\section{Verification purpose}

\textbf{Case ID:} \texttt{\CaseID}. This unit test verifies that ELMFIRE transports a circular level-set front at a prescribed, spatially uniform normal rate. It exercises raster ingestion, the \texttt{ADJ} spread-rate multiplier, the level-set update, and \texttt{phi}/\texttt{vs} output without wind, slope, crown fire, spotting, or suppression. It can reveal rate-unit errors, incorrect signs, and failure to couple the calculated rate to front motion. This is code verification against a known solution, not comparison with a physical fire. The verification question is whether the final spread-rate and radius satisfy the two stated tolerances at \(t=360\) s.

\section{Mathematical formulation and numerical implementation}

ELMFIRE advances the signed level set by
\[
 \frac{\partial\phi}{\partial t}+V_n|\nabla\phi|=0,
\]
with \(\phi<0\) burned and \(\phi=0\) the front. Uniform rasters and disabled directional forcing make \(V_n=R_0=\ReferenceROS\). The configured \texttt{SPREAD\_RATE\_IN\_M} option makes the \texttt{vs} output use \(\mathrm{m\,min^{-1}}\), while integration time remains seconds. Raster georeferencing supplies the mesh; no obsolete computational-domain namelist is used.

\section{Derivation of expected behavior}

For initial circle radius \(r_0=5\) m and constant outward speed \(R_0\), radial symmetry reduces the level-set equation to \(dr/dt=R_0\). Hence
\[
 r(360)=5+(0.0116332)(360)=9.187952~\mathrm{m},
 \qquad R_0=0.697992~\mathrm{m\,min^{-1}}.
\]
These are independent consequences of the prescribed input rate and elapsed time.

% BEGIN EXPLICIT METRIC DERIVATION
The numerical reference values and the reported metrics follow directly from the generated rate raster.  The raster value is \(R_s=0.0116332~\mathrm{m\,s^{-1}}\), so the report-unit spread rate is
\[
 R_{\mathrm{ref}}=60R_s=60(0.0116332)=0.697992~\mathrm{m\,min^{-1}}.
\]
With \(r_0=5~\mathrm{m}\) and \(T=360~\mathrm{s}\), the exact final radius is
\[
 r_{\mathrm{ref}}=r_0+R_sT=5+(0.0116332)(360)=9.187952~\mathrm{m}.
\]
The two calculated errors are \(e_R=|R_{\max}-0.697992|/0.697992\) and
\(e_r=|\sqrt{N_b(1~\mathrm{m})^2/\pi}-9.187952|\), where \(N_b\) is the
number of final cells with \(\phi<0\).  Thus the acceptance substitutions are
\(e_R\leq0.05\) and \(e_r\leq2~\mathrm{m}\).
% END EXPLICIT METRIC DERIVATION

\section{Verification metrics and acceptance criteria}

The maximum finite final \texttt{vs} value is compared with \(0.697992~\mathrm{m\,min^{-1}}\); a 5 percent relative tolerance catches conversion or propagation-rate errors while allowing front reconstruction effects. The equivalent burned radius
\(r_A=\sqrt{N_b\Delta x^2/\pi}\), where \(N_b\) counts valid final cells with \(\phi<0\), is compared with 9.187952 m. Area-equivalent radius is less sensitive to pixel stair-stepping than a single axial intercept; its tolerance is two cells. Nodata cells are excluded. The lexicographically last matching final \texttt{phi} and \texttt{vs} rasters are required. The case passes iff both components pass; missing or malformed required output is NOT EVALUABLE.

\section{Simulation configuration and predicted outcome}

The generated EPSG:32610 domain is 200 m square with 200 cells per side (\(\Delta x=1\) m), uniform fuel/moisture/terrain/weather fields, and a centered 5 m ignition circle. Duration is \FinalTime. The predicted final field is circular, with rate \(0.697992~\mathrm{m\,min^{-1}}\) and equivalent radius 9.187952 m.

\begin{center}\begin{tabular}{p{0.22\textwidth}p{0.22\textwidth}p{0.12\textwidth}p{0.34\textwidth}}\toprule
Parameter & Configured value & Units & Role \\
\midrule Domain/grid & 200 by 200 cells; 1 & m/cell & Resolves the manufactured circle \\
Initial front & centered radius 5 & m & Exact initial condition \\
Rate & 0.0116332 & m/s & Constant manufactured normal speed \\
Time control & step 1; stop 360 & s & Exact final time and timestep multiple \\
Disabled forcing & wind, slope, crown, spotting, suppression & --- & Isolates level-set transport \\
\bottomrule\end{tabular}\end{center}

Figure~\ref{fig:input-configuration} is generated mechanically from the prepared case-local GeoTIFF inputs. It identifies the selected source artifact and representative variant; the left panel shows a fuel, terrain, moisture, wind, canopy, or other controlling field selected for this experiment, and the right panel shows the cells selected by the non-positive initial level-set field. No values in this figure are redrawn or inferred from the expected response.

\begin{figure}[H]
  \centering
  \IfFileExists{../figures/input_configuration.pdf}{%
    \includegraphics[width=0.98\textwidth,height=0.46\textheight,keepaspectratio]{input_configuration.pdf}%
  }{\fbox{\parbox{0.88\textwidth}{The prepared-input figure has not been generated. Run the case-local preprocessing and spatial-evidence step.}}}
  \caption{Prepared input conditions from case-local generated rasters. Coordinates are in metres; the colorbar identifies the displayed field and its physical units or category. The ignition panel is derived directly from \(\phi\leq0\).}
  \label{fig:input-configuration}
\end{figure}

\section{Actual simulation results}

Only case-local ELMFIRE rasters are read. Current status is \textbf{\VerificationStatus}; \PassingMetricCount{} of \MetricCount{} components pass.
\begin{figure}[H]\centering
\IfFileExists{../figures/whole_domain_result.pdf}{\includegraphics[width=0.72\textwidth]{whole_domain_result.pdf}}{\fbox{No case-local whole-domain result is available.}}
\caption{Final whole-domain signed level-set field for the completed variant at the final output time, selected from the case-local data output directory. Values are dimensionless signed-level-set values; the generated table supplies the acceptance limits.}
\end{figure}

\section{Calculated metrics and verification decision}

\begin{center}
\small\setlength{\tabcolsep}{4pt}
\begin{tabular}{@{}llll@{}}
\toprule
Metric & Expected value or limit & Calculated value & Status \\
\midrule
\MetricRows
\bottomrule
\end{tabular}
\end{center}
The authoritative overall decision is \textbf{\VerificationStatus}. The generated table is sourced directly from the case-local metrics JSON and is not recomputed in this report.

\section{Technical interpretation}

A spread-rate failure points to input scaling, units, or \texttt{vs} selection. A radius-only failure indicates level-set integration or front/area discretization. A missing raster is an evidence failure, not a scientific pass.

\section{Scope and limitations}

This case verifies one deterministic constant-rate, flat, no-wind propagation state. It does not validate a fuel model, quantify anisotropy, or test coupled crown, spotting, or suppression behavior.

